% GENERATED FILE. Edit canonical lessons, then run npm run generate:books.
% canonical-source-hash: fnv1a64:5e7378d5df34dea8
% canonical-lessons: AR-C513-sukkari, AR-C513-rabw, AR-C513-iaqa, AR-C513-daufa, AR-C513-nazafa

\chapter{Long illnesses}
\label{ch:ar-long-illnesses}
\hlchaptermodality{513}

\begin{chapteropening}
\textbf{By the end of this chapter:} \emph{I can talk about diabetes, asthma and a disability, and say that someone grew weak or bled.}

Complete the last lesson of chapter 513: I can talk about diabetes, asthma and a disability, and say that someone grew weak or bled.
\end{chapteropening}

\section[\texorpdfstring{\ar{سكري} (sukkarī)}{سكري (sukkarī)}]{\ar{سكري} (sukkarī) — diabetes}

\label{lesson:AR-C513-sukkari}



\begin{quote}
Before the new one: say the Arabic for good health, then the Arabic for to be in pain.
\end{quote}

\subsection*{You'll want to know: \ar{سكري}}
\textbf{\ar{سكري}} — \emph{sukkarī} — \textquotedblleft{}diabetes\textquotedblright{}.

\begin{quote}
\ar{عند} \ar{جدي} \ar{مرض} \ar{السكري}.
\end{quote}

\begin{quote}
\emph{ʿinda jaddī maraḍ as-sukkarī.}
\end{quote}

\begin{quote}
My grandfather has diabetes.
\end{quote}

The full name is \ar{مرض} \ar{السكري}, the sugar illness.

\begin{cousinweb}
\textbf{\ar{سكري}} (\emph{sukkarī}) is built on \emph{sukkar}, \textquotedblleft{}sugar\textquotedblright{}, with the ending \emph{-ī}, \textquotedblleft{}of\textquotedblright{}: the sugar illness. \emph{Sukkar} itself is usually traced through Persian to Sanskrit \emph{śarkarā}.
\end{cousinweb}

\begin{grammarlens}[title={What you've built}]
One more word for describing what you lived and felt.
\end{grammarlens}

\subsection*{Your turn}
\begin{itemize}
\raggedright
  \item \emph{Say it:} \emph{sukkarī}
  \item \emph{Say it:} \emph{sukkarī}, once more
  \item \emph{Say it:} \emph{ʿinda jaddī maraḍ as-sukkarī.} — My grandfather has diabetes.
  \item \emph{Recall:} say the Arabic for good health, then the Arabic for to be in pain, then say \emph{sukkarī} again
\end{itemize}

\begin{tcolorbox}[breakable,colback=teal!4,colframe=teal!35!black,title={Before you move on}]
What does \ar{سكري} mean? (\textquotedblleft{}Diabetes\textquotedblright{}.) Say it once more.
\end{tcolorbox}
\section[\texorpdfstring{\ar{ربو} (rabw)}{ربو (rabw)}]{\ar{ربو} (rabw) — asthma}

\label{lesson:AR-C513-rabw}



\begin{quote}
Before the new one: say the Arabic for to be in pain, then the Arabic for diabetes.
\end{quote}

\subsection*{You'll want to know: \ar{ربو}}
\textbf{\ar{ربو}} — \emph{rabw} — \textquotedblleft{}asthma\textquotedblright{}.

\begin{quote}
\ar{يحمل} \ar{أخي} \ar{دواء} \ar{الربو} \ar{في} \ar{حقيبته}.
\end{quote}

\begin{quote}
\emph{yaḥmilu akhī dawāʾ ar-rabw fī ḥaqībatihi.}
\end{quote}

\begin{quote}
My brother carries his asthma medicine in his bag.
\end{quote}

\begin{cousinweb}
\textbf{\ar{ربو}} (\emph{rabw}) is built on the native Arabic root \emph{r-b-w}, \textquotedblleft{}to grow; to be out of breath\textquotedblright{} (Lane). Lane gives rabw as being out of breath, panting after running.
\end{cousinweb}

\begin{grammarlens}[title={What you've built}]
One more word for describing what you lived and felt.
\end{grammarlens}

\subsection*{Your turn}
\begin{itemize}
\raggedright
  \item \emph{Say it:} \emph{rabw}
  \item \emph{Say it:} \emph{rabw}, once more
  \item \emph{Say it:} \emph{yaḥmilu akhī dawāʾ ar-rabw fī ḥaqībatihi.} — My brother carries his asthma medicine in his bag.
  \item \emph{Recall:} say the Arabic for to be in pain, then the Arabic for diabetes, then say \emph{rabw} again
\end{itemize}

\begin{tcolorbox}[breakable,colback=teal!4,colframe=teal!35!black,title={Before you move on}]
What does \ar{ربو} mean? (\textquotedblleft{}Asthma\textquotedblright{}.) Say it once more.
\end{tcolorbox}
\section[\texorpdfstring{\ar{إعاقة} (iʿāqa)}{إعاقة (iʿāqa)}]{\ar{إعاقة} (iʿāqa) — a disability}

\label{lesson:AR-C513-iaqa}



\begin{quote}
Before the new one: say the Arabic for diabetes, then the Arabic for asthma.
\end{quote}

\subsection*{You'll want to know: \ar{إعاقة}}
\textbf{\ar{إعاقة}} — \emph{iʿāqa} — \textquotedblleft{}a disability\textquotedblright{}.

\begin{quote}
\ar{يعمل} \ar{في} \ar{الشركة} \ar{موظف} \ar{عنده} \ar{إعاقة}.
\end{quote}

\begin{quote}
\emph{yaʿmalu fī ash-sharika muwaẓẓaf ʿindahu iʿāqa.}
\end{quote}

\begin{quote}
An employee with a disability works at the company.
\end{quote}

\begin{cousinweb}
\textbf{\ar{إعاقة}} (\emph{iʿāqa}) is built on the native Arabic root \emph{ʿ-w-q}, \textquotedblleft{}to hinder, to hold back\textquotedblright{} (Lane). Iʿāqa is the verbal noun of Form IV: a hindrance.
\end{cousinweb}

\begin{grammarlens}[title={What you've built}]
One more word for describing what you lived and felt.
\end{grammarlens}

\subsection*{Your turn}
\begin{itemize}
\raggedright
  \item \emph{Say it:} \emph{iʿāqa}
  \item \emph{Say it:} \emph{iʿāqa}, once more
  \item \emph{Say it:} \emph{yaʿmalu fī ash-sharika muwaẓẓaf ʿindahu iʿāqa.} — An employee with a disability works at the company.
  \item \emph{Recall:} say the Arabic for diabetes, then the Arabic for asthma, then say \emph{iʿāqa} again
\end{itemize}

\begin{tcolorbox}[breakable,colback=teal!4,colframe=teal!35!black,title={Before you move on}]
What does \ar{إعاقة} mean? (\textquotedblleft{}A disability\textquotedblright{}.) Say it once more.
\end{tcolorbox}
\section[\texorpdfstring{\ar{ضعف} (ḍaʿufa)}{ضعف (ḍaʿufa)}]{\ar{ضعف} (ḍaʿufa) — to grow weak}

\label{lesson:AR-C513-daufa}



\begin{quote}
Before the new one: say the Arabic for asthma, then the Arabic for a disability.
\end{quote}

\subsection*{You'll want to know: \ar{ضعف}}
\textbf{\ar{ضعف}} — \emph{ḍaʿufa} — \textquotedblleft{}to grow weak\textquotedblright{}.

\begin{quote}
\ar{ضعف} \ar{جدي} \ar{بعد} \ar{المرض}.
\end{quote}

\begin{quote}
\emph{ḍaʿufa jaddī baʿda al-maraḍ.}
\end{quote}

\begin{quote}
My grandfather grew weak after his illness.
\end{quote}

Present: \ar{يضعف}, \emph{yaḍʿuf}.

\begin{cousinweb}
\textbf{\ar{ضعف}} (\emph{ḍaʿufa}) is built on the native Arabic root \emph{ḍ-ʿ-f}, \textquotedblleft{}to be weak\textquotedblright{} (Lane). Ḍaʿīf, weak, which you know, is from the same root.
\end{cousinweb}

\begin{grammarlens}[title={What you've built}]
One more word for describing what you lived and felt.
\end{grammarlens}

\subsection*{Your turn}
\begin{itemize}
\raggedright
  \item \emph{Say it:} \emph{ḍaʿufa}
  \item \emph{Say it:} \emph{ḍaʿufa}, once more
  \item \emph{Say it:} \emph{ḍaʿufa jaddī baʿda al-maraḍ.} — My grandfather grew weak after his illness.
  \item \emph{Recall:} say the Arabic for asthma, then the Arabic for a disability, then say \emph{ḍaʿufa} again
\end{itemize}

\begin{tcolorbox}[breakable,colback=teal!4,colframe=teal!35!black,title={Before you move on}]
What does \ar{ضعف} mean? (\textquotedblleft{}To grow weak\textquotedblright{}.) Say it once more.
\end{tcolorbox}
\section[\texorpdfstring{\ar{نزف} (nazafa)}{نزف (nazafa)}]{\ar{نزف} (nazafa) — to bleed}

\label{lesson:AR-C513-nazafa}



\begin{quote}
Before the new one: say the Arabic for a disability, then the Arabic for to grow weak.
\end{quote}

\subsection*{You'll want to know: \ar{نزف}}
\textbf{\ar{نزف}} — \emph{nazafa} — \textquotedblleft{}to bleed\textquotedblright{}.

\begin{quote}
\ar{نزف} \ar{أنف} \ar{اللاعب} \ar{بعد} \ar{الضربة}.
\end{quote}

\begin{quote}
\emph{nazafa anf al-lāʿib baʿda aḍ-ḍarba.}
\end{quote}

\begin{quote}
The player's nose bled after the blow.
\end{quote}

Present: \ar{ينزف}, \emph{yanzif}.

The noun \ar{نزيف}, bleeding, shares the root.

\begin{cousinweb}
\textbf{\ar{نزف}} (\emph{nazafa}) is built on the native Arabic root \emph{n-z-f}, \textquotedblleft{}to drain, to exhaust\textquotedblright{} (Lane). Lane gives nazafa as to drain a well dry; bleeding drains the blood.
\end{cousinweb}

\begin{grammarlens}[title={What you've built}]
One more word for describing what you lived and felt.
\end{grammarlens}

\subsection*{Your turn}
\begin{itemize}
\raggedright
  \item \emph{Say it:} \emph{nazafa}
  \item \emph{Say it:} \emph{nazafa}, once more
  \item \emph{Say it:} \emph{nazafa anf al-lāʿib baʿda aḍ-ḍarba.} — The player's nose bled after the blow.
  \item \emph{Recall:} say the Arabic for a disability, then the Arabic for to grow weak, then say \emph{nazafa} again
\end{itemize}

\begin{tcolorbox}[breakable,colback=teal!4,colframe=teal!35!black,title={Before you move on}]
What does \ar{نزف} mean? (\textquotedblleft{}To bleed\textquotedblright{}.) Say it once more.
\end{tcolorbox}
